\chapter{Math}
\label{chap:math}

The \texttt{std::math} module provides mathematical operations optimized for 68k Amigas. Since most classic Amigas lack a floating-point unit (FPU), the library emphasizes fixed-point arithmetic, lookup tables, and integer math---all designed for predictable performance without FPU dependencies.

\section{Core Integer Math}
\label{sec:math-core}

The \texttt{std::math::core} module provides fundamental integer operations with optimized 68020+ assembly implementations.

\subsection{Importing Core Math}

\begin{lstlisting}
from std::math::core import (
    abs, min, max, clamp, sign, diff,
    is_power_of_two, next_power_of_two,
    div_ceil, div_round
)
\end{lstlisting}

\subsection{Absolute Value}

The \texttt{abs()} function returns the absolute value of a signed integer:

\begin{lstlisting}
let x: i32 = -42
let positive = abs(x)  // Returns 42

let small: i16 = -100
let pos_small = abs(small)  // Returns 100
\end{lstlisting}

The function is overloaded for \texttt{i8}, \texttt{i16}, and \texttt{i32} types. The implementation uses optimized assembly that avoids branches on 68020+.

\subsection{Minimum and Maximum}

Find the smaller or larger of two values:

\begin{lstlisting}
let a: i32 = 10
let b: i32 = 25

let smaller = min(a, b)  // Returns 10
let larger = max(a, b)   // Returns 25

// Works with unsigned types too
let x: u32 = 100
let y: u32 = 50
let result = max(x, y)  // Returns 100
\end{lstlisting}

\begin{notebox}
\texttt{min()} and \texttt{max()} are overloaded for all integer types: \texttt{i8}, \texttt{i16}, \texttt{i32}, \texttt{u8}, \texttt{u16}, and \texttt{u32}.
\end{notebox}

\subsection{Clamping Values}

The \texttt{clamp()} function restricts a value to a range:

\begin{lstlisting}
let value: i32 = 150
let clamped = clamp(value, 0, 100)  // Returns 100

let low_value: i32 = -50
let clamped_low = clamp(low_value, 0, 100)  // Returns 0

let in_range: i32 = 50
let unchanged = clamp(in_range, 0, 100)  // Returns 50
\end{lstlisting}

This is essential for keeping coordinates on screen, volumes within range, or any bounded value.

\subsection{Sign Function}

The \texttt{sign()} function returns -1, 0, or 1 based on the value's sign:

\begin{lstlisting}
sign(-42)  // Returns -1
sign(0)    // Returns 0
sign(100)  // Returns 1
\end{lstlisting}

\subsection{Absolute Difference}

The \texttt{diff()} function returns the absolute difference between two values:

\begin{lstlisting}
diff(10, 25)   // Returns 15
diff(25, 10)   // Returns 15 (order doesn't matter)
diff(-5, 5)    // Returns 10
\end{lstlisting}

\subsection{Power of Two Operations}

Test whether a value is a power of two, or find the next power of two:

\begin{lstlisting}
is_power_of_two(16)   // true
is_power_of_two(17)   // false
is_power_of_two(256)  // true
is_power_of_two(0)    // false

next_power_of_two(5)    // Returns 8
next_power_of_two(16)   // Returns 16 (already power of 2)
next_power_of_two(100)  // Returns 128
\end{lstlisting}

Power-of-two sizes are often required for Amiga hardware resources like bitplanes and copper lists.

\subsection{Division with Rounding}

Standard integer division truncates toward zero. These functions provide alternative rounding:

\begin{lstlisting}
// Round up (ceiling division)
div_ceil(10, 3)   // Returns 4 (10/3 = 3.33, rounds up)
div_ceil(9, 3)    // Returns 3 (exact division)

// Round to nearest
div_round(10, 3)  // Returns 3 (10/3 = 3.33, rounds to nearest)
div_round(11, 3)  // Returns 4 (11/3 = 3.67, rounds to nearest)
\end{lstlisting}

\section{Fixed-Point Arithmetic}
\label{sec:math-fixed}

For 68k Amigas without FPUs, fixed-point math provides fractional precision using integer operations. The \texttt{std::math::fixed} module provides two fixed-point types optimized for different use cases.

\subsection{Understanding Fixed-Point}

Fixed-point numbers store fractional values using integers by dedicating some bits to the fractional part. For example, a 16.16 format uses 16 bits for the integer part and 16 bits for the fraction.

The value stored is: \texttt{actual\_value = raw\_value / scale\_factor}

\begin{comingfromc}
In C, fixed-point math requires manual scaling and tracking. Novus provides typed wrappers that handle scaling automatically while generating the same efficient code.
\end{comingfromc}

\subsection{Fixed32: 16.16 Format}

\texttt{Fixed32} uses a 16.16 format (16 integer bits, 16 fractional bits):

\begin{lstlisting}
from std::math::fixed import Fixed32

// Create from integer
let ten = Fixed32::from_i32(10)

// Create from integer and fractional parts
let mixed = Fixed32::from_parts(3, 32768)  // 3.5 (32768 = 0.5 * 65536)

// Create from raw value (already scaled)
let raw = Fixed32::from_raw(65536)  // 1.0 (65536 = 1 << 16)
let half = Fixed32::from_raw(32768) // 0.5
\end{lstlisting}

\begin{table}[h]
\centering
\begin{tabular}{ll}
\textbf{Property} & \textbf{Value} \\
\hline
Format & 16.16 (signed) \\
Range & -32768.0 to +32767.99998 \\
Precision & ~0.000015 (1/65536) \\
Scale factor & 65536 \\
\end{tabular}
\caption{Fixed32 characteristics}
\end{table}

\subsubsection{Arithmetic Operations}

\begin{lstlisting}
let a = Fixed32::from_i32(10)
let b = Fixed32::from_parts(2, 32768)  // 2.5

// Basic arithmetic
let sum = a.add(b)        // 12.5
let diff = a.sub(b)       // 7.5
let product = a.mul(b)    // 25.0
let quotient = a.div(b)   // 4.0
let negated = a.neg()     // -10.0
let absolute = negated.abs()  // 10.0
\end{lstlisting}

\subsubsection{Rounding and Conversion}

\begin{lstlisting}
let value = Fixed32::from_parts(3, 32768)  // 3.5

let floored = value.floor()   // Fixed32 representing 3.0
let ceiled = value.ceil()     // Fixed32 representing 4.0
let rounded = value.round()   // Fixed32 representing 4.0

// Convert to integer
let int_floor = value.to_i32()          // 3 (truncates toward zero)
let int_round = value.to_i32_rounded()  // 4 (rounds to nearest)
\end{lstlisting}

\subsubsection{Comparison}

\begin{lstlisting}
let x = Fixed32::from_i32(5)
let y = Fixed32::from_parts(5, 32768)  // 5.5

x.lt(y)  // true (5 < 5.5)
x.le(y)  // true
x.gt(y)  // false
x.ge(y)  // false
x.eq(y)  // false
\end{lstlisting}

\subsubsection{Mathematical Constants}

\begin{lstlisting}
from std::math::fixed import (
    fixed32_pi, fixed32_half_pi,
    fixed32_two_pi, fixed32_sqrt2, fixed32_one
)

let pi = fixed32_pi()           // 3.14159...
let half_pi = fixed32_half_pi() // 1.57079...
let two_pi = fixed32_two_pi()   // 6.28318...
let sqrt2 = fixed32_sqrt2()     // 1.41421...
let one = fixed32_one()         // 1.0
\end{lstlisting}

\subsection{Fixed16: 8.8 Format}

For situations where memory is tight or less precision is acceptable, \texttt{Fixed16} uses an 8.8 format:

\begin{lstlisting}
from std::math::fixed import Fixed16

let x = Fixed16::from_i8(10)
let y = Fixed16::from_raw(64)  // 0.25 (64 = 0.25 * 256)

let sum = x.add(y)  // 10.25
\end{lstlisting}

\begin{table}[h]
\centering
\begin{tabular}{ll}
\textbf{Property} & \textbf{Value} \\
\hline
Format & 8.8 (signed) \\
Range & -128.0 to +127.996 \\
Precision & ~0.004 (1/256) \\
Scale factor & 256 \\
\end{tabular}
\caption{Fixed16 characteristics}
\end{table}

\texttt{Fixed16} is useful for:
\begin{itemize}
\item Sprite animation deltas (sub-pixel movement)
\item Color component interpolation
\item Compact storage when precision isn't critical
\end{itemize}

\subsection{Converting Between Fixed Types}

\begin{lstlisting}
from std::math::fixed import fixed32_to_fixed16, fixed16_to_fixed32

// Fixed32 to Fixed16 (loses precision)
let big = Fixed32::from_i32(5)
let small = fixed32_to_fixed16(big)

// Fixed16 to Fixed32 (preserves value)
let compact = Fixed16::from_i8(5)
let expanded = fixed16_to_fixed32(compact)
\end{lstlisting}

\begin{warningbox}
Converting from \texttt{Fixed32} to \texttt{Fixed16} can overflow if the value exceeds \texttt{Fixed16}'s range (-128 to +127).
\end{warningbox}

\section{Trigonometry}
\label{sec:math-trig}

The \texttt{std::math::trig} module provides sine, cosine, and related functions using lookup tables. This is essential for games and demos on FPU-less Amigas.

\subsection{Binary Angular Measurement (BAMS)}

Angles in Novus use \textbf{Binary Angular Measurement} (BAMS), where a full circle is represented by the range 0-65535 (\texttt{u16}):

\begin{table}[h]
\centering
\begin{tabular}{ll}
\textbf{BAMS Value} & \textbf{Degrees} \\
\hline
0 & 0\textdegree \\
16384 & 90\textdegree \\
32768 & 180\textdegree \\
49152 & 270\textdegree \\
65535 & ~360\textdegree \\
\end{tabular}
\caption{BAMS angle values}
\end{table}

BAMS has several advantages:
\begin{itemize}
\item Angles wrap naturally with 16-bit overflow (65535 + 1 = 0)
\item No need to normalize angles to 0-360 range
\item Fast lookup table indexing
\item Angle addition/subtraction works directly
\end{itemize}

\subsection{Sine and Cosine}

\begin{lstlisting}
from std::math::trig import sin_fixed32, cos_fixed32

// 90 degrees in BAMS
let angle: u16 = 16384

let s = sin_fixed32(angle)  // ~1.0
let c = cos_fixed32(angle)  // ~0.0

// Use for movement
let speed = Fixed32::from_i32(5)
let dx = speed.mul(cos_fixed32(angle))
let dy = speed.mul(sin_fixed32(angle))
\end{lstlisting}

\subsubsection{Raw Integer Results}

For performance-critical code, get raw 14-bit scaled values:

\begin{lstlisting}
from std::math::trig import sin_i16, cos_i16

let angle: u16 = 16384
let s = sin_i16(angle)  // 16384 (represents 1.0)
let c = cos_i16(angle)  // 0

// Scale: 16384 = 1.0, -16384 = -1.0
\end{lstlisting}

This avoids \texttt{Fixed32} wrapper overhead for blitter or copper calculations.

\subsubsection{Higher-Precision with Interpolation}

For smoother results, use interpolated variants:

\begin{lstlisting}
from std::math::trig import sin_fixed32_interp, cos_fixed32_interp

let angle: u16 = 1000  // Arbitrary angle
let s = sin_fixed32_interp(angle)  // Linear interpolation between table entries
\end{lstlisting}

Interpolated functions are slightly slower but provide smoother results for fine angle increments.

\subsection{Angle Conversion}

\begin{lstlisting}
from std::math::trig import (
    degrees_to_angle, angle_to_degrees,
    radians_to_angle, angle_to_radians
)
from std::math::fixed import fixed32_pi

// Degrees to BAMS
let angle = degrees_to_angle(45)  // Returns 8192

// BAMS to degrees
let deg = angle_to_degrees(16384)  // Returns 90

// Radians (as Fixed32) to BAMS
let pi = fixed32_pi()
let angle_rad = radians_to_angle(pi)  // Returns 32768 (180 degrees)
\end{lstlisting}

\subsection{Tangent}

\begin{lstlisting}
from std::math::trig import tan_fixed32

let angle: u16 = 8192  // 45 degrees
let t = tan_fixed32(angle)  // ~1.0
\end{lstlisting}

\begin{warningbox}
Tangent approaches infinity at 90\textdegree{} and 270\textdegree{}. The function returns clamped maximum values at these points.
\end{warningbox}

\subsection{Atan2: Angle from Coordinates}

Calculate the angle from x,y coordinates (essential for aiming, pathfinding):

\begin{lstlisting}
from std::math::trig import atan2_angle

// Point at (10, 10) from origin
let angle = atan2_angle(10, 10)  // ~8192 (45 degrees)

// Point at (-10, 0) from origin
let angle2 = atan2_angle(0, -10)  // 32768 (180 degrees)

// Use for "look at" behavior
let dx = target_x - entity_x
let dy = target_y - entity_y
let facing_angle = atan2_angle(dy, dx)
\end{lstlisting}

\subsection{Angle Constants}

\begin{lstlisting}
from std::math::trig import (
    ANGLE_FULL,           // 0 (360 degrees wraps to 0)
    ANGLE_HALF,           // 32768 (180 degrees)
    ANGLE_QUARTER,        // 16384 (90 degrees)
    ANGLE_THREE_QUARTER,  // 49152 (270 degrees)
    ANGLE_ONE_DEGREE      // 182 (~1 degree)
)
\end{lstlisting}

\section{Square Root and Distance}
\label{sec:math-sqrt}

The \texttt{std::math::sqrt} module provides square root operations with automatic FPU detection on equipped machines.

\subsection{Integer Square Root}

\begin{lstlisting}
from std::math::sqrt import sqrt

let root = sqrt(16u32)   // Returns 4
let root2 = sqrt(100u32) // Returns 10
let root3 = sqrt(2u32)   // Returns 1 (integer truncation)
\end{lstlisting}

The function is overloaded for \texttt{u32}, \texttt{i32}, and \texttt{u16}. For signed values, negative inputs return 0.

\subsection{Fixed-Point Square Root}

\begin{lstlisting}
let value = Fixed32::from_i32(2)
let root = sqrt(value)  // ~1.414 (sqrt of 2)
\end{lstlisting}

\begin{tipbox}
On systems with an FPU (68881, 68882, 68040, 68060), the compiler automatically uses hardware \texttt{FSQRT} instructions for dramatically faster execution.
\end{tipbox}

\subsection{Fast Approximate Square Root}

When precision is less important than speed:

\begin{lstlisting}
from std::math::sqrt import sqrt_fast

let approx = sqrt_fast(value)  // Newton-Raphson approximation
\end{lstlisting}

\subsection{Distance Calculations}

\begin{lstlisting}
from std::math::sqrt import distance, distance_less_than, distance_within

// Euclidean distance between two points
let d = distance(x1, y1, x2, y2)

// Optimized distance comparisons (avoid sqrt when possible)
if distance_less_than(x1, y1, x2, y2, 100) {
    // Points are less than 100 units apart
}

if distance_within(x1, y1, x2, y2, 100) {
    // Points are within 100 units (inclusive)
}
\end{lstlisting}

\subsection{Hypotenuse}

Overflow-safe hypotenuse calculation:

\begin{lstlisting}
from std::math::sqrt import hypot

// Calculate sqrt(x^2 + y^2) without overflow
let length = hypot(30000, 40000)  // Returns 50000
\end{lstlisting}

\subsection{Inverse Square Root}

For vector normalization:

\begin{lstlisting}
from std::math::sqrt import inv_sqrt

let value = Fixed32::from_i32(4)
let inv = inv_sqrt(value)  // 0.5 (1/sqrt(4))
\end{lstlisting}

\section{Angle Utilities}
\label{sec:math-angle}

The \texttt{std::math::angle} module extends trigonometry with angle manipulation functions.

\subsection{Angle Normalization}

\begin{lstlisting}
from std::math::angle import angle_normalize

// BAMS angles wrap automatically, but this ensures consistency
let normalized = angle_normalize(angle)
\end{lstlisting}

\subsection{Angle Difference}

Find the shortest angular distance between two angles:

\begin{lstlisting}
from std::math::angle import angle_diff, angle_diff_abs

let a: u16 = 1000   // Small angle
let b: u16 = 64000  // Near full circle

// Signed difference (shortest path)
let diff = angle_diff(a, b)  // Negative (b is "behind" a)

// Absolute difference
let abs_diff = angle_diff_abs(a, b)  // Always positive
\end{lstlisting}

This correctly handles the wraparound at 360\textdegree.

\subsection{Angle Interpolation}

Smoothly rotate between angles:

\begin{lstlisting}
from std::math::angle import angle_lerp, angle_smoothstep, angle_damp

let start: u16 = 0
let end: u16 = 16384  // 90 degrees

// Linear interpolation (t = 0.5)
let mid = angle_lerp(start, end, Fixed32::from_raw(32768))  // 45 degrees

// Smooth interpolation (ease in/out)
let smooth = angle_smoothstep(start, end, t)

// Frame-rate independent smooth rotation
let current = angle_damp(current, target, smoothness, delta_time)
\end{lstlisting}

\subsection{Angle Snapping}

Snap angles to fixed increments:

\begin{lstlisting}
from std::math::angle import angle_snap, angle_snap_45, angle_snap_8way

// Snap to arbitrary increment
let snapped = angle_snap(angle, 8192)  // Snap to 45-degree increments

// Convenience for common cases
let snapped45 = angle_snap_45(angle)   // Same as above
let snapped8 = angle_snap_8way(angle)  // 8 directions
\end{lstlisting}

\subsection{Direction Conversion}

Convert angles to discrete direction indices:

\begin{lstlisting}
from std::math::angle import angle_to_8way, angle_to_4way

// Get direction index (0-7) for sprite animation frames
let dir = angle_to_8way(facing_angle)

// Get cardinal direction index (0-3)
let cardinal = angle_to_4way(facing_angle)

// Direction values:
// 4-way: 0=right, 1=down, 2=left, 3=up
// 8-way: 0=right, 1=down-right, 2=down, 3=down-left, etc.
\end{lstlisting}

\subsection{Quadrant and Facing Detection}

\begin{lstlisting}
from std::math::angle import (
    quadrant, facing_right, facing_left, facing_up, facing_down
)

let angle: u16 = 8192  // 45 degrees

let q = quadrant(angle)  // Returns 0 (first quadrant)

// Boolean facing checks
if facing_right(angle) {
    // Angle is between -90 and +90 degrees
}

if facing_up(angle) {
    // Angle is between 0 and 180 degrees (upper half)
}
\end{lstlisting}

\section{Interpolation}
\label{sec:math-interp}

The \texttt{std::math::interp} module provides functions for smooth value transitions.

\subsection{Linear Interpolation (Lerp)}

\begin{lstlisting}
from std::math::interp import lerp, lerp_i32, lerp_clamped

let start = Fixed32::from_i32(0)
let end = Fixed32::from_i32(100)
let t = Fixed32::from_raw(16384)  // 0.25 (16384 = 0.25 * 65536)

let result = lerp(start, end, t)  // 25.0

// Integer version
let int_result = lerp_i32(0, 100, t)  // 25

// Clamped version (t is clamped to 0-1)
let safe = lerp_clamped(start, end, Fixed32::from_i32(2))  // 100 (not 200)
\end{lstlisting}

\subsection{Inverse Lerp and Remap}

\begin{lstlisting}
from std::math::interp import inv_lerp, remap

// Find t value for a result
let t = inv_lerp(0, 100, 25)  // 0.25

// Remap from one range to another
let remapped = remap(25, 0, 100, 0, 1000)  // 250
\end{lstlisting}

\subsection{Smoothstep}

Hermite interpolation with smooth acceleration/deceleration:

\begin{lstlisting}
from std::math::interp import smoothstep, smootherstep

let t = Fixed32::from_raw(32768)  // 0.5

// Standard smoothstep (3t^2 - 2t^3)
let smooth = smoothstep(t)

// Ken Perlin's smoother version (6t^5 - 15t^4 + 10t^3)
let smoother = smootherstep(t)
\end{lstlisting}

\subsection{Movement Functions}

\begin{lstlisting}
from std::math::interp import move_towards, damp, spring_damp

// Move toward target at fixed speed
let new_pos = move_towards(current, target, max_delta)

// Exponential decay (frame-rate independent smoothing)
let smoothed = damp(current, target, smoothness, delta_time)

// Spring physics (overshoots and settles)
let spring_pos = spring_damp(current, target, velocity, stiffness, damping, delta_time)
\end{lstlisting}

\subsection{Bezier Curves}

\begin{lstlisting}
from std::math::interp import bezier_quadratic, bezier_cubic

let t = Fixed32::from_raw(32768)  // 0.5

// Quadratic Bezier (3 control points)
let p0 = Fixed32::from_i32(0)
let p1 = Fixed32::from_i32(50)   // Control point
let p2 = Fixed32::from_i32(100)
let quad = bezier_quadratic(p0, p1, p2, t)

// Cubic Bezier (4 control points)
let c1 = Fixed32::from_i32(25)   // Control point 1
let c2 = Fixed32::from_i32(75)   // Control point 2
let cubic = bezier_cubic(p0, c1, c2, p2, t)
\end{lstlisting}

\subsection{Looping Functions}

\begin{lstlisting}
from std::math::interp import ping_pong, repeat

// Oscillate between 0 and length
let pp = ping_pong(time, length)  // 0 -> length -> 0 -> length...

// Repeat in range (like modulo but for Fixed32)
let r = repeat(value, length)  // 0 -> length, wraps to 0
\end{lstlisting}

\section{Easing Functions}
\label{sec:math-ease}

The \texttt{std::math::ease} module provides Robert Penner's easing equations for smooth animations.

\subsection{Easing Overview}

Each easing function takes a \texttt{t} value from 0.0 to 1.0 and returns a modified value. There are three variants for each curve:
\begin{itemize}
\item \texttt{ease\_in\_*}: Starts slow, accelerates
\item \texttt{ease\_out\_*}: Starts fast, decelerates
\item \texttt{ease\_in\_out\_*}: Smooth start and end
\end{itemize}

\subsection{Basic Usage}

\begin{lstlisting}
from std::math::ease import ease_in_quad, ease_out_quad, ease_in_out_quad

let t = Fixed32::from_raw(32768)  // 0.5

let slow_start = ease_in_quad(t)      // Accelerating
let slow_end = ease_out_quad(t)       // Decelerating
let smooth = ease_in_out_quad(t)      // Both
\end{lstlisting}

\subsection{Available Easing Curves}

\begin{table}[h]
\centering
\begin{tabular}{lp{8cm}}
\textbf{Family} & \textbf{Description} \\
\hline
\texttt{linear} & No easing (straight line) \\
\texttt{quad} & Quadratic ($t^2$) - gentle curve \\
\texttt{cubic} & Cubic ($t^3$) - more pronounced \\
\texttt{quart} & Quartic ($t^4$) - even stronger \\
\texttt{quint} & Quintic ($t^5$) - very strong \\
\texttt{sine} & Sinusoidal - very smooth \\
\texttt{expo} & Exponential - dramatic acceleration \\
\texttt{circ} & Circular - based on circle equation \\
\texttt{back} & Overshoots slightly, then settles \\
\texttt{elastic} & Spring-like bounce \\
\texttt{bounce} & Bouncing ball effect \\
\end{tabular}
\caption{Easing function families}
\end{table}

\subsection{Example: Animating Position}

\begin{lstlisting}
from std::math::ease import ease_out_cubic
from std::math::interp import lerp

struct Animation {
    start_x: i32,
    end_x: i32,
    duration: i32,     // Frames
    current_frame: i32,
}

impl Animation {
    fn update(&var self) -> i32 {
        if self.current_frame >= self.duration {
            return self.end_x
        }

        // Calculate t (0.0 to 1.0)
        // t = current_frame / duration (as fixed-point)
        let t = Fixed32::from_i32(self.current_frame)
            .div(Fixed32::from_i32(self.duration))

        // Apply easing
        let eased_t = ease_out_cubic(t)

        // Interpolate position
        self.current_frame++
        return lerp_i32(self.start_x, self.end_x, eased_t)
    }
}
\end{lstlisting}

\subsection{Special Easing Functions}

\subsubsection{Back (Overshoot)}

\begin{lstlisting}
from std::math::ease import ease_in_back, ease_out_back, ease_in_out_back

// Pulls back before moving forward
let backed = ease_in_back(t)
\end{lstlisting}

Good for UI elements that need a "windup" feel.

\subsubsection{Elastic (Spring)}

\begin{lstlisting}
from std::math::ease import ease_in_elastic, ease_out_elastic, ease_in_out_elastic

// Spring oscillation
let elastic = ease_out_elastic(t)
\end{lstlisting}

Creates a spring-like wobble effect.

\subsubsection{Bounce}

\begin{lstlisting}
from std::math::ease import ease_in_bounce, ease_out_bounce, ease_in_out_bounce

// Bouncing ball effect
let bounced = ease_out_bounce(t)
\end{lstlisting}

Simulates an object bouncing and settling.

\subsubsection{Steps}

\begin{lstlisting}
from std::math::ease import ease_steps

// Staircase animation (for pixel art, retro effects)
let stepped = ease_steps(t, 4)  // 4 discrete steps
\end{lstlisting}

\section{2D Vectors}
\label{sec:math-vec2}

The \texttt{std::math::vec2} module provides 2D vector types for games and graphics.

\subsection{Vec2: Fixed-Point Vectors}

\texttt{Vec2} uses \texttt{Fixed32} components for sub-pixel precision:

\begin{lstlisting}
from std::math::vec2 import Vec2

// Create vectors
let pos = Vec2::new(
    Fixed32::from_i32(100),
    Fixed32::from_i32(50)
)
let zero = Vec2::zero()
let unit_x = Vec2::unit_x()  // (1, 0)
let unit_y = Vec2::unit_y()  // (0, 1)

// From integers
let from_int = Vec2::from_i32(10, 20)
\end{lstlisting}

\subsection{Vec2i: Integer Vectors}

\texttt{Vec2i} uses \texttt{i32} components for pixel-perfect operations:

\begin{lstlisting}
from std::math::vec2 import Vec2i

let screen_pos = Vec2i::new(320, 200)
let origin = Vec2i::zero()
\end{lstlisting}

\subsection{Vector Operations}

\begin{lstlisting}
let a = Vec2::from_i32(10, 20)
let b = Vec2::from_i32(5, 10)

// Basic arithmetic
let sum = a.add(b)        // (15, 30)
let diff = a.sub(b)       // (5, 10)
let neg = a.neg()         // (-10, -20)

// Scaling
let scaled = a.scale(Fixed32::from_i32(2))  // (20, 40)
let halved = a.div(Fixed32::from_i32(2))    // (5, 10)

// Component-wise multiplication (Hadamard product)
let hadamard = a.hadamard(b)  // (50, 200)
\end{lstlisting}

\subsection{Dot and Cross Products}

\begin{lstlisting}
let a = Vec2::from_i32(1, 0)
let b = Vec2::from_i32(0, 1)

// Dot product (scalar)
let dot = a.dot(b)  // 0 (perpendicular)

// 2D cross product (scalar, represents z-component)
let cross = a.cross(b)  // 1 (counter-clockwise)
\end{lstlisting}

\subsection{Length and Distance}

\begin{lstlisting}
let v = Vec2::from_i32(3, 4)

// Squared length (fast, avoids sqrt)
let len_sq = v.length_sq()  // 25

// Approximate length (alpha-max-beta-min algorithm)
let len_approx = v.length_approx()  // ~5 (fast approximation)

// Distance between points
let a = Vec2::from_i32(0, 0)
let b = Vec2::from_i32(3, 4)
let dist_sq = distance_sq(a, b)        // 25
let dist_approx = distance_approx(a, b) // ~5
\end{lstlisting}

\begin{tipbox}
The alpha-max-beta-min algorithm provides ~3\% accuracy for length calculations without expensive square roots. It's ideal for collision detection where exact values aren't critical.
\end{tipbox}

\subsection{Perpendicular Vectors}

\begin{lstlisting}
let v = Vec2::from_i32(1, 0)

// Counter-clockwise perpendicular
let perp = v.perp()     // (0, 1)

// Clockwise perpendicular
let perp_cw = v.perp_cw()  // (0, -1)
\end{lstlisting}

\subsection{Reflection}

\begin{lstlisting}
let velocity = Vec2::from_i32(1, -1)  // Moving down-right
let normal = Vec2::from_i32(0, 1)     // Horizontal surface

let reflected = velocity.reflect(normal)  // Bounces up-right
\end{lstlisting}

\subsection{Rotation}

\begin{lstlisting}
from std::math::vec2 import Vec2
from std::math::trig import ANGLE_QUARTER

let v = Vec2::from_i32(1, 0)  // Pointing right

// Rotate 90 degrees counter-clockwise
let rotated = v.rotate(ANGLE_QUARTER)  // Now pointing up
\end{lstlisting}

\subsection{Direction and Angle}

\begin{lstlisting}
from std::math::vec2 import direction_from_angle, angle_from_direction

// Create unit vector from angle
let angle: u16 = 8192  // 45 degrees
let dir = direction_from_angle(angle)  // (~0.707, ~0.707)

// Get angle from vector
let v = Vec2::from_i32(1, 1)
let a = angle_from_direction(v)  // ~8192 (45 degrees)
\end{lstlisting}

\subsection{Interpolation}

\begin{lstlisting}
from std::math::vec2 import lerp_vec2

let start = Vec2::from_i32(0, 0)
let end = Vec2::from_i32(100, 50)
let t = Fixed32::from_raw(32768)  // 0.5

let mid = lerp_vec2(start, end, t)  // (50, 25)
\end{lstlisting}

\subsection{Converting Between Vec2 and Vec2i}

\begin{lstlisting}
let fixed_pos = Vec2::from_i32(10, 20)

// Convert to integer tuple (for screen drawing)
let (screen_x, screen_y) = fixed_pos.to_i32()  // Truncates fractional parts

// Create Vec2i from integers
let screen_pos = Vec2i::new(screen_x, screen_y)

// Convert Vec2i back to Vec2 (for physics)
let physics_pos = screen_pos.to_vec2()
\end{lstlisting}

\section{Practical Example: 2D Game Physics}
\label{sec:math-example}

Here's a complete example showing fixed-point math, vectors, and trigonometry working together:

\begin{lstlisting}
from std::math::fixed import Fixed32
from std::math::vec2 import Vec2, Vec2i, lerp_vec2
from std::math::trig import sin_fixed32, cos_fixed32, atan2_angle
from std::math::interp import move_towards
from std::math::core import clamp

struct Entity {
    pos: Vec2,           // Sub-pixel position
    velocity: Vec2,      // Movement per frame
    facing: u16,         // BAMS angle
    speed: Fixed32,
    max_speed: Fixed32,
}

impl Entity {
    fn new(x: i32, y: i32) -> Entity {
        return Entity {
            pos: Vec2::from_i32(x, y),
            velocity: Vec2::zero(),
            facing: 0,
            speed: Fixed32::from_i32(0),
            max_speed: Fixed32::from_i32(4),
        }
    }

    fn accelerate(&var self, angle: u16, accel: Fixed32) {
        // Calculate acceleration vector
        let ax = accel.mul(cos_fixed32(angle))
        let ay = accel.mul(sin_fixed32(angle))

        // Add to velocity
        self.velocity = self.velocity.add(Vec2::new(ax, ay))

        // Clamp to max speed
        let speed_sq = self.velocity.length_sq()
        let max_sq = self.max_speed.mul(self.max_speed)
        if speed_sq.gt(max_sq) {
            // Normalize and scale to max
            let len = self.velocity.length_approx()
            self.velocity = self.velocity.scale(self.max_speed.div(len))
        }

        // Update facing angle (threshold = 0.01)
        if speed_sq.gt(Fixed32::from_raw(655)) {  // ~0.01
            self.facing = atan2_angle(
                self.velocity.y.to_i32(),
                self.velocity.x.to_i32()
            )
        }
    }

    fn apply_friction(&var self, friction: Fixed32) {
        self.velocity = self.velocity.scale(friction)
    }

    fn update(&var self) {
        // Apply velocity to position
        self.pos = self.pos.add(self.velocity)
    }

    fn screen_pos(&self) -> Vec2i {
        // Convert fixed-point position to screen coordinates
        let (x, y) = self.pos.to_i32()
        return Vec2i::new(x, y)
    }
}

fn main() -> i32 {
    var player = Entity::new(160, 100)

    // Game loop simulation
    while var frame < 60 {
        // Accelerate right
        let right: u16 = 0
        player.accelerate(right, Fixed32::from_raw(16384))  // 0.25

        // Apply friction (0.95)
        player.apply_friction(Fixed32::from_raw(62259))  // ~0.95

        // Update position
        player.update()

        // Get screen position for drawing
        let screen = player.screen_pos()

        frame++
    }

    return 0
}
\end{lstlisting}

\section{Performance Considerations}
\label{sec:math-perf}

\subsection{Avoiding Expensive Operations}

\begin{table}[h]
\centering
\begin{tabular}{lll}
\textbf{Operation} & \textbf{Cost} & \textbf{Alternative} \\
\hline
Division & High & Multiply by reciprocal, bit shift \\
Square root & Very high & Use squared distances \\
Trig functions & Medium & Lookup tables (built-in) \\
Floating point & N/A & Fixed-point (no FPU) \\
\end{tabular}
\caption{Operation costs on 68k}
\end{table}

\subsection{Squared Distance Pattern}

\begin{lstlisting}
// Slow: calculates actual distance
if distance(a, b) < threshold {
    // ...
}

// Fast: compare squared values
let threshold_sq = threshold * threshold
if distance_sq(a, b) < threshold_sq {
    // ...
}
\end{lstlisting}

\subsection{Precomputing Constants}

\begin{lstlisting}
// Calculate once at startup
let inv_screen_width = Fixed32::from_i32(1).div(Fixed32::from_i32(320))

// Use in loop (multiply is faster than divide)
let normalized_x = x.mul(inv_screen_width)
\end{lstlisting}

\section{Processor-Specific Tuning}
\label{sec:math-processor}

The 68k processor family spans multiple generations, from the 68020 to the FPGA-based 68080. Novus targets 68020 and later processors, which provide native 32-bit multiply and divide operations essential for efficient fixed-point math. Understanding the differences between processors is essential for writing efficient code.

\subsection{Processor Overview}

\begin{table}[h]
\centering
\small
\begin{tabular}{lcccc}
\textbf{Processor} & \textbf{Multiply} & \textbf{Divide} & \textbf{FPU} & \textbf{Cache} \\
\hline
68020 & 32$\times$32 & 64$\div$32 & External & 256B I \\
68030 & 32$\times$32 & 64$\div$32 & External & 256B I+D \\
68040 & 32$\times$32 & 64$\div$32 & Built-in* & 4KB I+D \\
68060 & 32$\times$32 & 64$\div$32 & Built-in* & 8KB I+D \\
68080 & 32$\times$32$\rightarrow$64 & 64$\div$32 & Built-in & 32KB I+D \\
\end{tabular}
\caption{68k processor capabilities. *Some FPU instructions emulated in software.}
\end{table}

\subsection{68020/68030: The Baseline}

The 68020 (A1200 stock, some A3000) and 68030 (A3000, A4000, accelerators) add crucial features:

\begin{itemize}
\item \textbf{32-bit multiply}: \texttt{MULS.L}/\texttt{MULU.L} are 32$\times$32$\rightarrow$32 or 64
\item \textbf{64-bit divide}: \texttt{DIVS.L}/\texttt{DIVU.L} are 64$\div$32$\rightarrow$32
\item \textbf{Bit field operations}: \texttt{BFEXTU}, \texttt{BFINS}, etc.
\item \textbf{Instruction cache}: 256 bytes (68020) or 256+256 (68030)
\item \textbf{Optional FPU}: 68881 or 68882 coprocessor
\end{itemize}

\begin{table}[h]
\centering
\begin{tabular}{lrl}
\textbf{Operation} & \textbf{Cycles} & \textbf{Notes} \\
\hline
\texttt{MULU.L} & 28--44 & 32$\times$32$\rightarrow$32 \\
\texttt{MULU.L} (64-bit) & 28--44 & 32$\times$32$\rightarrow$64 \\
\texttt{DIVU.L} & 42--90 & 64$\div$32$\rightarrow$32 \\
\texttt{DIVS.L} & 58--122 & Signed division slower \\
\end{tabular}
\caption{68020/030 instruction timing (approximate cycles)}
\end{table}

\textbf{68020/030 Optimization Tips:}
\begin{itemize}
\item Use full 32-bit \texttt{Fixed32} operations---hardware supports them natively
\item Bit field instructions are powerful for packed data (sprites, graphics)
\item Loop unrolling helps due to small instruction cache
\item Consider 68881/68882 FPU if available (adds hardware sin, cos, sqrt, etc.)
\end{itemize}

\subsection{68040: Built-in FPU with Caveats}

The 68040 (A4000, many accelerators) integrates the FPU but with important limitations:

\begin{itemize}
\item \textbf{Partial FPU}: Some 68882 instructions are emulated in software
\item \textbf{Larger caches}: 4KB instruction + 4KB data
\item \textbf{Pipelined}: Better instruction throughput
\item \textbf{Copyback cache}: Can cause issues with DMA if not flushed
\end{itemize}

\begin{warningbox}
The 68040 FPU does \textbf{not} implement these 68882 instructions in hardware:
\texttt{FSIN}, \texttt{FCOS}, \texttt{FTAN}, \texttt{FASIN}, \texttt{FACOS}, \texttt{FATAN}, \texttt{FSINH}, \texttt{FCOSH}, \texttt{FTANH}, \texttt{FATANH}, \texttt{FLOG10}, \texttt{FLOG2}, \texttt{FLOGN}, \texttt{FMOD}, \texttt{FREM}, \texttt{FETOX}, \texttt{FGETEXP}, \texttt{FGETMAN}, \texttt{FTWOTOX}, \texttt{FTENTOX}.

These generate F-line exceptions and are handled by the \texttt{68040.library} FPSP (Floating Point Support Package). Software emulation is 10--100$\times$ slower than hardware.
\end{warningbox}

\textbf{Hardware FPU Operations (fast):}
\begin{itemize}
\item \texttt{FADD}, \texttt{FSUB}, \texttt{FMUL}, \texttt{FDIV}
\item \texttt{FSQRT} (hardware!)
\item \texttt{FABS}, \texttt{FNEG}, \texttt{FMOVE}
\item \texttt{FCMP}, \texttt{FTST}
\end{itemize}

\textbf{68040 Optimization Tips:}
\begin{itemize}
\item Use lookup tables for trig functions---faster than emulated \texttt{FSIN}/\texttt{FCOS}
\item \texttt{FSQRT} is hardware and fast---use it!
\item Flush caches before DMA operations
\item The math library uses hardware \texttt{FSQRT} automatically on 68040+
\end{itemize}

\subsection{68060: Fastest Classic 68k}

The 68060 (Blizzard, Apollo, etc.) is the pinnacle of classic 68k:

\begin{itemize}
\item \textbf{Superscalar}: Can execute 2 instructions per cycle
\item \textbf{Branch prediction}: Reduces pipeline stalls
\item \textbf{Faster multiply}: Single-cycle for common cases
\item \textbf{8KB caches}: Instruction and data
\item \textbf{Even more emulated FPU ops}: More instructions trap than 68040
\end{itemize}

\begin{table}[h]
\centering
\begin{tabular}{lrl}
\textbf{Operation} & \textbf{Cycles} & \textbf{Notes} \\
\hline
\texttt{MULU.L} & 2 & Most cases! \\
\texttt{MULS.L} & 2--4 & Depends on operands \\
\texttt{DIVU.L} & 35--37 & Still expensive \\
\texttt{FSQRT} & 1 & Hardware, pipelined \\
\texttt{FMUL} & 1 & Hardware, pipelined \\
\end{tabular}
\caption{68060 instruction timing (cycles)}
\end{table}

\begin{warningbox}
The 68060 emulates even more instructions than 68040. Additionally, these \textbf{integer} instructions trap and are emulated:
\begin{itemize}
\item \texttt{MOVEP} (move peripheral)
\item \texttt{CAS2} (compare and swap)
\item Some \texttt{CHK2}/\texttt{CMP2} forms
\item \texttt{MULU.L}/\texttt{MULS.L} with 64-bit result (Dh:Dl form)
\item \texttt{DIVU.L}/\texttt{DIVS.L} with 64-bit dividend (Dr:Dq form)
\end{itemize}

Always install the \texttt{68060.library} for proper emulation support.
\end{warningbox}

\textbf{68060 Optimization Tips:}
\begin{itemize}
\item Native multiply is so fast that shift/add sequences are often \textbf{slower}
\item Division is still expensive---precompute reciprocals
\item Use lookup tables for trig (emulated FPU trig is very slow)
\item Hardware \texttt{FSQRT} and \texttt{FMUL} are excellent
\item Align loops to cache lines for best branch prediction
\end{itemize}

\subsection{68080 (Apollo Core): Modern Performance}

The FPGA-based 68080 brings modern features to the 68k architecture:

\begin{itemize}
\item \textbf{AMMX SIMD}: 64-bit vector operations (8$\times$8-bit, 4$\times$16-bit, 2$\times$32-bit)
\item \textbf{Fast multiply}: 32$\times$32$\rightarrow$64 hardware multiply, fully pipelined
\item \textbf{Complete FPU}: All instructions in hardware, no emulation
\item \textbf{Large caches}: 32KB instruction + 32KB data
\item \textbf{Fast memory}: Direct SDRAM access
\end{itemize}

\begin{table}[h]
\centering
\begin{tabular}{lrl}
\textbf{Operation} & \textbf{Cycles} & \textbf{Notes} \\
\hline
\texttt{MULU.L} & 2--4 & Pipelined, depends on data \\
\texttt{DIVU.L} & 10--20 & Hardware but still iterative \\
\texttt{FSIN}/\texttt{FCOS} & 15--30 & Hardware, no emulation needed \\
\texttt{FSQRT} & 8--15 & Hardware, iterative algorithm \\
AMMX vector op & 1--2 & Up to 8 operations in parallel \\
\end{tabular}
\caption{68080 instruction timing (approximate)}
\end{table}

\textbf{68080 AMMX Features:}
\begin{itemize}
\item 32 64-bit data registers (D0--D7, E0--E23)
\item SIMD arithmetic: 8$\times$8-bit or 4$\times$16-bit parallel operations
\item Butterfly operations for DSP (FFT, audio processing)
\item 3-operand instructions: \texttt{D = A op B}
\item Similar to SSE/AltiVec on x86/PowerPC
\end{itemize}

\begin{notebox}
Novus does not yet generate AMMX instructions automatically. For AMMX optimization, use inline assembly with the \texttt{asm \{\}} block. Future versions may add AMMX intrinsics.
\end{notebox}

\subsection{FPU Availability Summary}

\begin{table}[h]
\centering
\small
\begin{tabular}{lccccc}
\textbf{System} & \textbf{CPU} & \textbf{FPU} & \textbf{FSQRT} & \textbf{FSIN/FCOS} & \textbf{Recommendation} \\
\hline
A1200 stock & 68EC020 & None & -- & -- & Fixed-point + tables \\
A1200 + 68030 & 68030 & Optional & HW/Emu & HW/Emu & Check at runtime \\
A3000 & 68030 & Optional & HW/Emu & HW/Emu & Check at runtime \\
A4000/030 & 68030 & Optional & HW/Emu & HW/Emu & Check at runtime \\
A4000/040 & 68040 & Built-in & HW & Emulated & Use tables for trig \\
Accelerators & 68060 & Built-in & HW & Emulated & Use tables for trig \\
Vampire/Icedrake & 68080 & Built-in & HW & HW & Full FPU available \\
\end{tabular}
\caption{FPU availability by system}
\end{table}

\subsection{Choosing the Right Approach}

\begin{table}[h]
\centering
\small
\begin{tabular}{lp{5cm}p{5cm}}
\textbf{Target} & \textbf{Best Approach} & \textbf{Avoid} \\
\hline
68020/030 & Fixed32, lookup tables for trig & Assuming FPU exists \\
68040/060 & Fixed32 or FPU for basic math, tables for trig & \texttt{FSIN}/\texttt{FCOS} (emulated) \\
68080 & Full FPU, consider AMMX for heavy math & None---everything is fast \\
Multi-target & Fixed32 + tables (works everywhere) & FPU-only code paths \\
\end{tabular}
\caption{Recommended math strategies by target}
\end{table}

\subsection{Runtime CPU Detection}

Novus provides CPU detection to choose optimal code paths:

\begin{lstlisting}
from std::os::cpu import get_cpu_type, has_fpu, CpuType
from std::math::sqrt import fixed32_sqrt
from std::math::fixed import Fixed32

// Choose optimal sqrt implementation based on detected CPU
fn compute_sqrt(value: Fixed32) -> Fixed32 {
    match get_cpu_type() {
        CpuType::M68040 | CpuType::M68060 | CpuType::M68080 => {
            if has_fpu() {
                // Hardware FSQRT available
                let f = value.to_f32();
                Fixed32::from_f32(f.sqrt())
            } else {
                fixed32_sqrt(value)
            }
        },
        _ => {
            // 68020/030: Use lookup table + interpolation
            fixed32_sqrt(value)
        }
    }
}
\end{lstlisting}

\begin{tipbox}
The standard library math functions already perform this detection internally. You typically don't need to check CPU type yourself unless writing custom optimized routines.
\end{tipbox}

\subsection{Compiler Target Flags}

Use compiler flags to optimize for specific processors:

\begin{lstlisting}[style=bash]
# Target specific processor
novus compile --cpu=68020 game.novus    # A1200/A3000 minimum (default)
novus compile --cpu=68040 game.novus    # 68040 optimizations
novus compile --cpu=68060 game.novus    # 68060 optimizations

# Multi-target fat binary
novus compile --cpu=fat:68020,68040,68060 game.novus
\end{lstlisting}

With \texttt{--cpu=68060}, the compiler:
\begin{itemize}
\item Uses native multiply instead of shift/add sequences
\item Avoids instructions that trap (like \texttt{MOVEP})
\item Generates cache-aligned loop entry points
\item Uses \texttt{FSQRT} for square root operations (when FPU available)
\end{itemize}

\section{Module Reference}
\label{sec:math-reference}

\begin{table}[h]
\centering
\begin{tabular}{lp{8cm}}
\textbf{Module} & \textbf{Contents} \\
\hline
\texttt{std::math::core} & abs, min, max, clamp, sign, diff, power-of-two, division \\
\texttt{std::math::fixed} & Fixed32, Fixed16, constants \\
\texttt{std::math::trig} & sin, cos, tan, atan2, angle conversion \\
\texttt{std::math::sqrt} & Square root, distance, hypot, inv\_sqrt \\
\texttt{std::math::angle} & Angle diff, lerp, snap, direction conversion \\
\texttt{std::math::interp} & lerp, smoothstep, bezier, spring, ping\_pong \\
\texttt{std::math::ease} & Complete easing function library \\
\texttt{std::math::vec2} & Vec2 (fixed), Vec2i (integer), operations \\
\end{tabular}
\caption{Math module summary}
\end{table}
